\documentclass[11pt,a4paper]{article}
\usepackage[T1]{fontenc}
\usepackage{lmodern}
\usepackage[margin=24mm,headheight=14pt]{geometry}
\usepackage{amsmath,amssymb,amsthm,mathtools}
\usepackage{microtype,booktabs,array,tabularx,longtable}
\usepackage[dvipsnames]{xcolor}
\usepackage{xurl}
\usepackage[colorlinks=true,linkcolor=MidnightBlue,urlcolor=MidnightBlue,citecolor=MidnightBlue]{hyperref}
\usepackage{enumitem,listings,fancyhdr,needspace}
\setlist{nosep,leftmargin=*}
\setlength{\parskip}{4pt}
\setlength{\parindent}{0pt}
\setlength{\emergencystretch}{2em}
\setcounter{tocdepth}{1}
\widowpenalty=10000
\clubpenalty=10000
\displaywidowpenalty=10000
\postdisplaypenalty=1000
\pagestyle{fancy}
\fancyhf{}
\fancyhead[L]{\small A fixed Diophantine representation of $\mathrm{Con}(\mathrm{ZFC})$}
\fancyhead[R]{\small ProveIt research contribution}
\fancyfoot[C]{\thepage}
\newtheorem{theorem}{Theorem}[section]
\newtheorem{lemma}[theorem]{Lemma}
\newtheorem{corollary}[theorem]{Corollary}
\newtheorem{proposition}[theorem]{Proposition}
\theoremstyle{definition}\newtheorem{definition}[theorem]{Definition}
\newcommand{\N}{\mathbb N}
\newcommand{\Z}{\mathbb Z}
\newcommand{\ZFC}{\mathsf{ZFC}}
\newcommand{\Con}{\operatorname{Con}}
\newcommand{\pc}{\operatorname{popcount}}
\newcommand{\MF}{\overline M_F}
\newcommand{\coeff}{\mathcal C}
\newcommand{\file}[1]{\texttt{\detokenize{#1}}}
\newcommand{\repo}{\url{https://github.com/VladimirReshetnikov/ProveIt}}
\lstset{basicstyle=\small\ttfamily,breaklines=true,columns=fullflexible,keepspaces=true,frame=single,rulecolor=\color{black!20},backgroundcolor=\color{black!2},aboveskip=8pt,belowskip=8pt}
\hypersetup{pdftitle={A fixed Diophantine representation of Con(ZFC)},pdfauthor={Prepared for Vladimir Reshetnikov's ProveIt repository},pdfsubject={Seven equations, sixteen positive unknowns, exact fixed coefficients}}

\title{A fixed Diophantine representation of\texorpdfstring{\\$\Con(\ZFC)$}{ Con(ZFC)}}
\author{Prepared for Vladimir Reshetnikov's \textsc{ProveIt} repository\\\small Mathematical construction and reproducible arithmetic artifacts}
\date{22 October 2026\\\small Based on snapshot \texttt{cdd582bec13f908df3711fa9e05d1ddac368a5db}}

\begin{document}
\maketitle
\begin{abstract}
We give an explicit system of seven ordinary polynomial equations in sixteen
strictly positive integer unknowns whose lack of solutions is equivalent,
provably in ZFC, to $\Con(\ZFC)$. Its largest equation degree is 28; a shared
arithmetic circuit uses 63 binary operations. The same construction gives
a single polynomial of degree 101 in 74 operations, or a single polynomial
of degree 157 in 72 operations. A separate explicit lift gives 70 quadratic
equations in 79 nonnegative unknowns and a single quartic with the same
solvability predicate.

Every coefficient is fixed by a terminating finite recipe applied to the
literal K215 transition table. No unspecified universal-machine index is
left over. The coefficients are enormous and are specified exactly in
compressed form, not expanded into decimal numerals. The new contribution
is a mathematical proof and executable source audit, using the repository's
existing K215 calibration and Diophantine compiler theorems. It is not a
new end-to-end Lean or Rocq formalization. No global minimality is claimed.
\end{abstract}

\tableofcontents
\newpage

\section{The recommended system}
\label{sec:system}

All five coefficient ports below denote fixed integers:
\[
 B-1,\quad K,\quad M_C,\quad \MF,\quad W_0.
\]
Their complete specification is in Section~\ref{sec:constants}; in
particular,
\begin{equation}\label{eq:visibleconstants}
 b=5^{43},\qquad d=5^{87},\qquad B=2^d,\qquad
 W_0=2^{2d+b}.
\end{equation}
The other three coefficient ports encode a single explicitly specified
finite local rule set obtained from K215. They are not unknowns.

Take precisely sixteen strictly positive integer unknowns, in this order:
\begin{equation}\label{eq:variables}
 J,F,\alpha,t,f,h,i,T,s,w,\tau,\eta,\zeta,y,Z,\gamma.
\end{equation}
Define the following polynomial abbreviations. They add no existential
variables and use no variable exponentiation:
\begin{align}
 q&=(B-1)J+1,& X&=wq,&Y&=sq^3,\notag\\
 E&=XY,& k&=\eta+\zeta,&c&=kY+\eta,\notag\\
 a&=E+Y,&H&=4a+3,&\Delta&=a^2+H,\notag\\
 D&=X+ac+\gamma H,&C&=Z+W_0,&G&=q^2-qF-Z,\notag\\
 R&=G(q^2-1)+(M_C+q\MF)J,\notag\\
 V&=c(Tf-1)-Rf^2,&L&=E(kY),\notag\\
 I&=ic^2,&K_a&=\Delta(f^2-1).
 \label{eq:abbreviations}
\end{align}
The recommended system $\mathcal S$ is
\begin{equation}\label{eq:system}
\boxed{\begin{aligned}
 \tau^2-L(L+k)&=1,\\
 D^2-\Delta c^2&=1,\\
 K_a(V^2-y^2)+y^2&=1,\\
 k-hE-R&=1,\\
 (K+w)C+(q-F)-t(q-1)&=1,\\
 I^2&=K_a,\\
 q-F-Z-\alpha&=C.
\end{aligned}}
\end{equation}
For example, the sixth equation is exactly
$(ic^2)^2=\Delta(f^2-1)$. It does not impose $\Delta\mid i$.
The source file \file{fixed_input63_ordinary_system.json} gives every
arithmetic instruction and the seven comparisons. Its witness domain is
part of the specification: allowing arbitrary signed values in
\eqref{eq:variables} without a domain conversion is not the same problem.

\begin{theorem}[Fixed ZFC instance]\label{thm:main}
For the coefficients of Section~\ref{sec:constants}, ZFC proves
\[
 \bigl(\exists\boldsymbol{x}\in\N_{>0}^{16}\;
                  \mathcal S(\boldsymbol{x})\bigr)
 \quad\longleftrightarrow\quad
 \neg\Con(\ZFC).
\]
Consequently ZFC proves
\[
 \boxed{\quad
 \bigl(\nexists\boldsymbol{x}\in\N_{>0}^{16}\;
                  \mathcal S(\boldsymbol{x})\bigr)
 \quad\longleftrightarrow\quad \Con(\ZFC).
 \quad}
\]
Here $\Con(\ZFC)$ is the usual arithmetical assertion that no finite
ZFC derivation proves a contradiction, with the proof calculus calibrated
by the pinned K215 endpoint.
\end{theorem}

\subsection{Single-polynomial versions}
First define two parent certificates whose unit-sign proofs are especially
transparent. One final simplification below removes a gate from each.
Write
\begin{align*}
 N_0&=\tau^2-L(L+k),&N_1&=D^2-\Delta c^2,\\
 N_a&=K_a(V^2-y^2)+y^2,&N_k&=k-hE-R,\\
 N_t&=(K+w)C+(q-F)-t(q-1),&N_s&=1+I^2-K_a,\\
 M&=1+(q-F-Z-\alpha-C)^2.
\end{align*}
Then a complete single equation is
\begin{equation}\label{eq:p75}
 \boxed{P_{75}:=N_0N_1N_aN_kN_tN_sM-1=0.}
\end{equation}
It has exactly the same positive solution set as
\eqref{eq:system}. The reverse implication is not a generic consequence
of taking products: Section~\ref{sec:proof} proves that all possibly
negative unit factors have the required positive signs.

There is a slightly shorter circuit if higher degree is acceptable.
Keep the common abbreviations, but set
\[
 \widetilde K_a=(i\Delta c^2)^2,\quad
 \widetilde N_a=\widetilde K_a(V^2-y^2)+y^2,\quad
 \widetilde N_s=\Delta f^2-\widetilde K_a.
\]
The alternative equation is
\begin{equation}\label{eq:p73}
 \boxed{P_{73}:=
 N_0N_1\widetilde N_aN_kN_t\widetilde N_sM-\Delta=0.}
\end{equation}
It has a positive solution if and only if \eqref{eq:system} has one.
The solution sets are not asserted to coincide with the same $i$:
a solution of \eqref{eq:p73} gives an ordinary-strong solution by
$i_{\rm ordinary}=\Delta i_{\rm normalized}$. Soundness in the other
direction is proved directly, not by dividing an arbitrary $i$ by
$\Delta$.

For the final shortest circuits, put
\[
 r_0=q-F-Z-\alpha-C
\]
and remove the addition of 1 from the marker factor. The resulting
single equations are
\begin{align}
 \boxed{P_{74}:=N_0N_1N_aN_kN_tN_s\,r_0^2-1=0,}
 \label{eq:p74}\\
 \boxed{P_{72}:=N_0N_1\widetilde N_aN_kN_t\widetilde N_s\,r_0^2-\Delta=0.}
 \label{eq:p72}
\end{align}
Their positive solvability is equivalent to that of
\eqref{eq:system}. Their individual positive zero tuples need not
satisfy the old budget equality: now $r_0=\pm1$. The precise
soundness and completeness argument is in
Section~\ref{sec:squaremarker}. The seven-equation source remains
unchanged.

A direct sum-of-squares polynomial, often easiest to trust mechanically, is
\begin{equation}\label{eq:sos}
 P_{\rm SOS}=\sum_{j=1}^7 r_j^2,
\end{equation}
where $r_j$ is the left side minus the right side of the $j$th equation
in \eqref{eq:system}. Integer solvability of $P_{\rm SOS}=0$ with the
specified positive primary domain is exactly system solvability.

\Needspace{9\baselineskip}
\subsection{What ``simplest'' means here}
There is no single agreed ordering of variable count, degree, coefficient
size, expanded term count, formula length and shared arithmetic cost.
The seven-equation system is our recommendation because it is short to
read, has sixteen unknowns and exposes the proof obligations directly.
The alternatives make the tradeoffs explicit.

\begin{table}[ht]
\centering\small
\begin{tabular}{@{}lrrrr@{}}
\toprule
Presentation & Unknowns & Equations & Degree & Operations\\
\midrule
Recommended system $\mathcal S$ & $16$ positive & $7$ & $28$ max. & $63$\\
$P_{74}$ & $16$ positive & $1$ & $101$ & $74$\\
$P_{72}$ & $16$ positive & $1$ & $157$ & $72$\\
$P_{\rm SOS}$ & $16$ positive & $1$ & $56$ & $83$\\
Integer SOS & $64$ integers & $1$ & $112$ & $211$\\
Quadratic lift & $79$ nonnegative & $70$ & $2$ max. & ---\\
Quartic lift & $79$ nonnegative & $1$ & $4$ & ---\\
\bottomrule
\end{tabular}
\caption{Exact counts for the supplied presentations. The circuit count
charges each binary addition, subtraction and multiplication, including
one multiplication per square and arithmetic involving fixed numerals.
Reuse and comparisons are free. Coefficient bit lengths are not charged.}
\label{tab:tradeoffs}
\end{table}

The pinned repository also proves the existence of a nine-unknown
representation for every recursively enumerable set, through
\file{Jones1982.nine_unknowns} in \file{NineMain.lean}~\cite{nine}.
Specializing that theorem to the inconsistency predicate gives a
nine-unknown existence result. It does not make the displayed sixteen
variables minimal, nor does it supply a smaller literal coefficient
packet in this contribution. Our target is an inspectable concrete
equation tied to the literal K215 machine, with short source and checked
arithmetic transformations.

\section{Exact fixed coefficients}\label{sec:constants}

\subsection{The finite machine underlying every coefficient}
The source is the literal binary table
\path{Computability/BusyBeaver/Research/ZFC/candidate-k215/k215.tm}
at the pinned snapshot~\cite{k215}. The supplied copy has SHA-256
\begin{center}\small\ttfamily
fdbcef0825886035c3091c8537a69535bc96f9e265f0ae45863543d40ae70090.
\end{center}
It has 430 transition rows, initial state 0, working states 0 through 214,
halt state 215, and blank symbol 0 on a two-way tape. Its only transition
into the halt state is \file{68 1 1 R 215}.

The explicit program \file{normalized_k215.py} turns this into a
436-state, eight-symbol, one-sided machine $\mathcal M$ with a stationary
first step and absorbing acceptance. Physical cell $-j$ stores the
original cells $j$ and $-j-1$ as two bits. An origin flag marks cell 0.
The eight tape symbols encode
\[
 (0,0,0),(1,0,0),(0,0,1),(0,1,0),
 (1,1,0),(0,1,1),(1,0,1),(1,1,1)
\]
in that order. States 0 through 5 are START, AFTER, ERASE, RETURN, SINK
and HALT. State $6+2q+e$, for $0\le q<215$ and $e\in\{0,1\}$,
simulates original state $q$ on the nonnegative or negative half of
the original tape. All 3,488 transitions are explicitly defined.

An input of $n\ge3$ ones is erased, the head returns to the marked
origin, and blank K215 simulation begins after $2n+2$ initialization
steps. A virtual move across $0/-1$ switches the selected track while
staying at the physical origin. Other moves follow the folded tape.
The tape and head correspondence is an inductive invariant, proving
for every such input that $\mathcal M$ accepts if and only if K215
halts blank. This conclusion is an all-time simulation proof, not an
inference from the finite execution tests.

The finite window compiler uses raw input 2 when its arithmetic input
is fixed to $x=1$. The normalized machine therefore starts with four
ones. It deliberately erases its input, so this interface convention
does not affect the desired halting predicate. It would be incorrect
to claim that this particular machine recognizes only the even inputs.

\subsection{The ordered window list}
The tile alphabet has size
\[
 a_{\rm tile}=2+8(436+1)+4=3502.
\]
IDs 0 and 1 are the vertical and horizontal boundary tiles. They are
followed by the payload tiles $(\text{symbol},\text{head})$, in
lexicographic order with head values None, $0,\ldots,435$, followed
by the four initial phases L, I, Q, R. The exact local validity test is
the total finite predicate \file{local_valid} in the supplied source.
It enforces boundary propagation, the phase grammar
LL, LI, II, IQ, QR, RR, one-step tape updates, no escaping head, and
acceptance on a terminal row.

Place the following two row-major $3\times3$ windows first:
\begin{align*}
 w_0&=(1,1,1,3499,3499,3500,439,439,878),\\
 w_1&=(1,1,1,3498,3499,3499,2,439,439).
\end{align*}
List all other locally valid nine-tuples in lexicographic order. This
defines every $w_j$ without reference to a run length or a halting
oracle. The four corner coordinates do not affect the center predicate,
but are retained in each window and in the overlap constraints.

The exact five-position cross count is $136449173881235$, so the
number of windows is
\begin{equation}\label{eq:windowcount}
 k_{\rm win}=136449173881235\cdot3502^4
 =20522746353535702537070419760.
\end{equation}
The formula, its derivation from transition counts, and independent small
exhaustive counting controls are in \file{normalization.md} and
\file{window_count_checks.json}. The actual enormous ordered list has
not been expanded. Its direct iterator examines at most $3502^9$
candidates, a fixed finite bound.

\subsection{Clause layout}
For the rest of this section put $a_t=3502$ and $k_t=k_{\rm win}$.
Use distinct symbols for layout constants to avoid confusion with the
polynomial abbreviations of Section~\ref{sec:system}:
\begin{align*}
 A_t&=2^{95},&
 p_t&=(k_t-73736)/2,\\
 j_\mu&=31522+p_t=10261373176767851268535204534,\\
 \mu_t&=A_t-2+\sum_{j=1}^{j_\mu}A_t^j,\\
 m_t&=k_t+15a_t+1=20522746353535702537070472291.
\end{align*}
For $r,c\in\{0,1,2\}$ and $0\le s<a_t$ define
\[
 e(r,c,s)=k_t+3a_t+(2-r)3a_t+(2-c)a_t+s,
 \qquad e_*=k_t+12a_t.
\]
The native position set is
\[
 \mathcal E=\{0,\ldots,k_t-1\}
 \cup\{k_t+3a_t,\ldots,k_t+12a_t\}
 \cup\{m_t,3m_t,9m_t,27m_t\}.
\]
The final member of the middle interval is one ignored dummy position.
Further define
\begin{align*}
 h_t&=51m_t+3a_t+1,&t_1&=105m_t+4a_t+2,\\
 t_2&=159m_t+4a_t+3,&g_t&=186m_t+4a_t+4,\\
 \mathcal L&=5^{44},&b&=5^{43},\quad d=b\mathcal L=5^{87},\\
 \mathcal R&=2^b,&B&=2^d.
\end{align*}
These are the compiler's literal least-power-of-five choices, not
arbitrary upper bounds. In particular,
\[
 5^{43}<213m_t+4a_t+5\le5^{44},\qquad
 5^{42}<95j_\mu+2\le5^{43}.
\]
The detailed bit-length bound in \file{constants.md} proves that
$95j_\mu+2$ is the required radix bit length before taking the next
power of five.

Each selector window contributes the fixed clause coefficient
\[
 c_j=1+\sum_{r,c=0}^{2}
 A_t^{1+(3r+c)a_t+(w_j)_{3r+c}}
 +\sum_{v=0}^{3}A_t^{1+9a_t+v}
 \quad(0\le j<k_t).
\]
For the copy positions and anchors use
\[
 c_{e(r,c,s)}=A_t^{1+(3r+c)a_t+s},\qquad
 c_{\lambda_v m_t}=A_t^{1+9a_t+v},
 \quad(\lambda_0,\lambda_1,\lambda_2,\lambda_3)=(1,3,9,27).
\]
Set $c_{e_*}=0$. This defines the complete finite coefficient map.

\subsection{The five actual coefficient ports}
Define
\begin{align*}
 D_C^0={}&\mathcal R^{3a_t}+\mathcal R^{h_t+a_t}
       +\mathcal R^{8m_t}+\mathcal R^{24m_t}\\
 &+\sum_{e\in\mathcal E\setminus\{e_*\}}
 c_e\bigl(\mathcal R^{t_1-e}+\mathcal R^{t_2-e}\bigr),\\
 D_R={}&\mathcal R^{h_t},\\
 \epsilon_t={}&\begin{cases}
 1,&5\mid 2D_C^0-D_R,\\
 0,&5\nmid 2D_C^0-D_R,
 \end{cases}\\
 D_C={}&D_C^0+\epsilon_t\mathcal R^{g_t},\qquad K=D_C+BD_R.
\end{align*}
The Boolean $\epsilon_t$ is a fixed integer determined by a finite
expression, not an extra Diophantine variable. Its actual K215 value
has not been evaluated; the formula fixes it uniquely. The correction
ensures $2D_C-D_R\not\equiv0\pmod5$.

The masks are
\begin{align*}
 M_C={}&B-1-\sum_{e\in\mathcal E\setminus\{1\}}\mathcal R^e
             -2\mathcal R^{e_*},\\
 M_F^{\rm nat}={}&4+\mu_t(\mathcal R^{t_1}+\mathcal R^{t_2})\\
 &+\sum_{\substack{r=0,1;\ c=0,1,2\\0\le s<a_t}}
                  \mathcal R^{e(r,c,s)}
 +\sum_{\substack{r=0,1,2;\ c=0,1\\0\le s<a_t}}
                  \mathcal R^{h_t+e(r,c,s)}\\
 &+\mathcal R^{9m_t}+\mathcal R^{27m_t},\\
 \MF={}&M_F^{\rm nat}+B-1,\qquad
 W_0=\mathcal R B^2=2^{2d+b}.
\end{align*}
The native masks satisfy
\begin{gather*}
 0<M_C,M_F^{\rm nat}<B-1,\qquad
 M_C\equiv2\pmod4,\quad M_F^{\rm nat}\equiv4\pmod8,\\
 \pc(M_C)+\pc(M_F^{\rm nat})=d.
\end{gather*}
The polynomial source name \file{MF} denotes $\MF$, the \emph{shifted}
coefficient. Applying the native range bound to this shifted value
would be an error. The source aliases are summarized in
Table~\ref{tab:ports}.

\begin{table}[ht]
\centering
\begin{tabular}{@{}ll@{}}\toprule
Source port & Exact mathematical value\\\midrule
\file{Bm1} & $B-1$\\
\file{Kconstant} & $D_C+BD_R$\\
\file{MC} & $M_C$\\
\file{MF} & $M_F^{\rm nat}+B-1$\\
\file{W0} & $2^{2d+b}$\\\bottomrule
\end{tabular}
\caption{The five coefficient ports are fixed numerals, not parameters
over which the theorem quantifies.}\label{tab:ports}
\end{table}

This is an ordinary polynomial specification: every exponent in this
section is fixed independently of the sixteen unknowns. All coefficient
computations have fixed finite bounds. No predicate such as ``is a
proof'' or ``halts'' is hidden in a coefficient definition.

The coefficient representation is compact, but the integers are not
small. $B$ alone has $5^{87}+1$ binary digits, approximately
$6.46\times10^{60}$ bits. We do not claim to have materialized these
integers. \file{fixed_compiler.py}, \file{fixed_constants.json} and
\file{constants.md} provide the exact algorithm and symbolic values.

\section{Proof of the arithmetic equivalence}\label{sec:proof}

We prove that the fixed system and both compact single polynomials have
a positive solution exactly when $\mathcal M$ accepts raw input 2.
The proof uses the pinned fixed-program compiler theorem~\cite{raw75},
its canonical completeness construction~\cite{f84}, and the specific
Pell rank and step-down lemmas~\cite{rank,reverse,asymmetric}.
These are ordinary mathematical dependencies; the accompanying finite
tests do not replace them.

For $A\ge2$, define the Pell sequences by
\[
 \chi_A(n)+\psi_A(n)\sqrt{A^2-1}
       =(A+\sqrt{A^2-1})^n\quad(n\ge0).
\]
Both sequences have integer values and satisfy their usual second-order
recurrences. Every positive solution to
$d^2-(A^2-1)c^2=1$ is
$(d,c)=(\chi_A(p),\psi_A(p))$ for some $p\ge1$.

\subsection{Recovering the unit signs}
At a positive zero of $P_{75}$, its seven integer factors multiply to 1.
Each is therefore $1$ or $-1$. Since $M\ge1$, it is 1, giving
\begin{equation}\label{eq:budget}
 q=F+2Z+\alpha+W_0,\quad C=Z+W_0,\quad
 0<C<q,\quad F+Z<q,\quad W_0<q.
\end{equation}
Also $\Delta=(a+2)^2-1$ is 0 or 3 modulo 4. The main norm $N_1$
cannot be $-1$ modulo 4. The ordinary strong factor is
$N_s=1+I^2-\Delta(f^2-1)$. If $\Delta\equiv0$ it is
$1+I^2$ modulo 4; if $\Delta\equiv3$ it is $I^2+f^2$ modulo 4.
Neither case can give 3. Hence $N_s=1$ and $K_a=I^2>0$.
Now $N_a$ is congruent to $y^2$ or $V^2$ modulo 4, so $N_a=1$.

For the first factor put $V_0=XY^2>1$. Since $L=V_0k$,
\[
 N_0=\tau^2-V_0(V_0+1)k^2.
\]
The negative Pell equation for $V_0(V_0+1)$ has no integer solution.
Indeed, suppose $x^2-V_0(V_0+1)z^2=-1$ with the smallest positive $z$.
Then $V_0z<x<(V_0+\tfrac12)z$. Multiplication by the inverse
unit $(2V_0+1)-2\sqrt{V_0(V_0+1)}$ produces another solution with
\[
 0<z'=(2V_0+1)z-2x<z,
\]
a contradiction. Thus $N_0=1$.

Write the two remaining signs as
\[
 N_k=\varepsilon,\quad N_t=\nu,
 \qquad\varepsilon,\nu\in\{-1,1\},\quad\varepsilon\nu=1.
\]
Their positivity will be proved, rather than assumed from the product.

\subsection{Preliminary bounds and Pell indices}
The actual native mask bounds, the shifted-mask identity and
\eqref{eq:budget} give the pretyping estimates
\begin{equation}\label{eq:bounds}
 (2q-1)(q^2-1)<R<q^4-q^3,\quad
 E=wsq^4\ge q^4>R+2,\quad a>R+2,
\end{equation}
with $q\ge B\ge16$. They do not assume that $q$ is dyadic or that
a valid computation has already been decoded.

Put $A=a+2$ and $P_0=2XY^2+1$. The fundamental positive solution of
the first Pell equation has $k=2$, so
\[
 \tau=\chi_{P_0}(n),\qquad k=2\psi_{P_0}(n),\qquad n\ge1.
\]
Since $P_0\equiv1\pmod E$ and $k=R+\varepsilon+hE$,
\[
 2n\equiv R+\varepsilon\pmod E,\qquad
 n\ge(R-1)/2\ge24.
\]
The main norm gives
$D=\chi_A(p)$ and $c=\psi_A(p)$ for $p\ge1$.
As $P_0>A$ and $c>k$, Pell monotonicity implies $p>n$.
Elementary growth bounds now give
\begin{equation}\label{eq:bootstrap}
 c>A\Delta^2,\qquad c>2p,\qquad c>2R.
\end{equation}
These bounds precede every use of auxiliary positivity or strong rank.

The ordinary strong equality gives
\[
 f^2=1+i^2c^4/\Delta,\qquad
 f>c^2/\sqrt\Delta>c^2/A>2c.
\]
Thus $c>\Delta+1$, $f^2>\Delta+c$, and $R<a<\Delta$ imply
\begin{equation}\label{eq:auxgap}
 K_a-(Rf^2+c)=(\Delta-R)f^2-\Delta-c>0.
\end{equation}
Using the stronger normalized-source bound $f>c^2$ here would be
unjustified; \eqref{eq:bootstrap} is what makes the ordinary version work.

\subsection{Restoring positive auxiliary witnesses}
The auxiliary norm is
\[
 K_aV^2-(K_a-1)y^2=1.
\]
If $v=|V|>1$, then
$v^2-1=(K_a-1)(y^2-v^2)$, whence $y\ge v+1$ and
$v\ge2K_a-1$. But $T>0$ and \eqref{eq:auxgap} give
\[
 V=c(Tf-1)-Rf^2>-Rf^2-c>-K_a.
\]
This excludes negative $V$ with $|V|>1$. The cases $V=0,1,-1$
are excluded by the norm or by $V\equiv-c\pmod f$ and
$2<c<f/2$. Therefore $V>1$ and $y>V$.

The generic strong-rank lemma~\cite{rank} applies to
$c=\psi_A(p)>A\Delta^2$ and $(ic^2)^2=\Delta(f^2-1)$.
It supplies an index $m$ with
\begin{equation}\label{eq:rank}
 f=\chi_A(m),\quad p\mid m,\quad c\mid m,\quad
 ic^2=\Delta\psi_A(m),\quad m\ge c>2p.
\end{equation}
In particular,
\[
 f^2-1=\Delta\psi_A(m)^2=ic^2\psi_A(m),
\]
so $c^2\mid f^2-1$. This integrality fact is essential.
Define
\begin{equation}\label{eq:restoredaux}
 o=cT-Rf,\qquad j=Tf-1-\frac{R(f^2-1)}c.
\end{equation}
They are integers and satisfy $of-c=V$ and $jc=V+R$.
Since the right sides are positive, $o,j>0$.

The signed auxiliary step-down lemma~\cite{reverse} now has precisely
its required hypotheses: \eqref{eq:rank},
$V=of-c=jc-R>0$, the auxiliary norm, and $0<p,R<c/2$.
It yields $p=R$. Its proof classifies the auxiliary Pell solution at
an odd index, uses its congruences modulo $f$ and $c$, then excludes
the negative and wrapped target alternatives using these strict bounds.

Write $2n=R+\varepsilon+vE$. Because $0<R+\varepsilon<E$, $v\ge0$.
If $v\ge1$, then $2n\ge R-1+E>2R+1$, contradicting $n<p=R$.
Therefore $2n=R+\varepsilon$. If $\varepsilon=-1$ then $p=2n+1$.
Pell duplication and monotonicity give
\[
 c=\psi_A(2n+1)>\psi_A(2n)\ge A k>k(Y+1),
\]
contrary to $kY<c<k(Y+1)$, which follows from $\eta,\zeta>0$.
Hence $\varepsilon=\nu=1$. Every equation of
\eqref{eq:system} has now been recovered from $P_{75}=0$.
The converse follows by direct multiplication.

\subsection{Recovering the power and the omitted fixed-input bridge}
Now $R=2n-1$ is odd. Write $R=2r+1$. The input-independent
asymmetric Pell-ratio argument~\cite{asymmetric} gives
\[
 Y>\frac{(X+1)^{2r}}{2X^r}-1>\frac{X^r}{3},\qquad
 a>\frac{X^{r+1}}3>2^R.
\]
For $E_A(v)=\chi_A(v)-a\psi_A(v)$, the initial values are 1 and 2,
and its recurrence gives $E_A(v)\equiv2^v\pmod H$.
The main root says $E_A(R)=X+\gamma H$. Since both $X$ and $2^R$
are strictly between 0 and $a<H$, their congruence forces
\begin{equation}\label{eq:power}
 X=2^R.
\end{equation}
Because $X=wq$, $q$ is a power of two, say $q=2^\ell$.
The lower bound in \eqref{eq:bounds} gives $R>3q>3\ell$,
so $q^3\mid X$ and $w_{\rm raw}=X/q^3>0$.
Also $(B-1)\mid(q-1)$ implies $q=B^N$ for some $N\ge1$.

Set $u=2d+b$, so $W_0=2^u$. The budget yields $u<\ell<R$.
The fixed-input replacement omitted the old subtraction $2d$ from
the budget; its positive restoration must be justified. Since
$q$ and $W_0$ are distinct powers of two, $q\ge2W_0$, and
\begin{equation}\label{eq:rawslack}
 q-C=\frac{q-W_0+F+\alpha}{2}>\frac{W_0}{2}\ge2d.
\end{equation}
Consequently $\alpha_{\rm raw}=q-C-2d>0$ is a valid raw slack.
The tempting formula $\alpha-2d$ need not be positive and is not used.

Restore the deleted input variables by
\begin{align*}
 \kappa&=\psi_A(u),&\mu&=\chi_A(u),\\
 \delta&=(\kappa-u)/\Delta,&\phi&=c-\kappa,\\
 \rho&=(E_A(u)-W_0)/H.
\end{align*}
All five are positive integers. Since $u$ is odd,
$\psi_A(u)\equiv u\pmod\Delta$, and growth gives
$\psi_A(u)>u$. As $u<R$, $\phi>0$.
The projected recurrence proves $\rho$ integral, and
$E_A(u)>2^u$ proves $\rho>0$. They satisfy exactly
\begin{align*}
 \kappa&=u+\delta\Delta,&c&=\kappa+\phi,\\
 \mu^2&=1+\Delta\kappa^2,&
 \mu&=W_0+a\kappa+\rho H.
\end{align*}
Finally let $z_{\rm raw}=t-1+wC>0$. The transport equation gives
\[
 (K+X)C-F=(t-1+wC)(q-1).
\]
Together with \eqref{eq:restoredaux}, $i_{\rm raw}=i$,
$w_{\rm raw}$ and $\alpha_{\rm raw}$, these are all nineteen
equations of the original raw75 certificate at arithmetic input 1.
Every restored coordinate is positive. Its fixed-program compiler
theorem therefore proves acceptance of raw input 2.

\subsection{Completeness and the 73-operation alternative}
If $\mathcal M$ accepts raw input 2, the pinned canonical F84
construction~\cite{f84} provides its eighteen positive coordinates at
arithmetic input 1. Its input bridge has $W=W_0$.
Discard $\delta,\rho,\sigma$, set
\[
 \gamma=\rho+\sigma,\qquad
 \alpha_{\rm new}=\alpha_{\rm F84}+2d,
\]
and keep the other coordinates. The resulting sixteen-tuple satisfies
the normalized equations and $P_{73}=0$. In that polynomial,
\[
 \widetilde N_s=\Delta
       \bigl(f^2-\Delta i^2c^4\bigr),\qquad \Delta>0.
\]
Canceling $\Delta$ turns its product equation into a product of integer
units. The same sign and index argument works, with the simpler strong
bound $f^2=1+\Delta i^2c^4$. Its full proof is supplied in
\file{FIXED_INPUT73_PROOF.md}.

Map this normalized solution to the ordinary one by
$i_{\rm ordinary}=\Delta i_{\rm normalized}$. Then
\[
 \Delta(f^2-1)=(i_{\rm ordinary}c^2)^2
              =(i_{\rm normalized}\Delta c^2)^2,
\]
so the auxiliary coefficient and all other required equations agree.
This proves completeness of \eqref{eq:system} and $P_{75}$.
Both compact polynomials therefore represent the same fixed
acceptance predicate.

\subsection{The final squared-marker saving}\label{sec:squaremarker}
At a zero of $P_{74}$, or of $P_{72}$ after canceling $\Delta>0$,
the square factor $r_0^2$ is an integer unit. Hence $r_0=\epsilon$
with $\epsilon\in\{-1,1\}$. The new budget is
\[
 q=W_0+F+2Z+\alpha+\epsilon.
\]
Since $\alpha+\epsilon\ge0$, it still implies $C<q$, $F+Z<q$
and $W_0<q$. These are all the budget consequences used before
dyadic power recovery. The norm signs, first and main Pell indices,
auxiliary positivity, rank, signed target recovery and no-wrap
argument therefore remain valid.

Once $q$ is known to be a power of two, $q\ge2W_0$ and
\[
 q-C=\frac{q-W_0+F+\alpha+\epsilon}{2}
      >\frac{W_0}{2}\ge2d,
\]
because $F+\alpha+\epsilon\ge1$. Thus the original strictly
positive slack $\alpha_{\rm raw}=q-C-2d$ is still restored.
All other restorations are unchanged. This proves soundness directly
through the full raw75 compiler, including the case
$\alpha=1$, $\epsilon=-1$.

For completeness, take a positive zero of $P_{75}$ or $P_{73}$
and replace $\alpha$ by $\alpha+1$. Its old budget residual was zero;
its new residual is $-1$, with square 1. No other factor uses
$\alpha$, so this gives a zero of $P_{74}$ or $P_{72}$.
The equivalence concerns existence of positive zeros. A pointwise
inverse using $\alpha+\epsilon$ would incorrectly allow a zero
coordinate in the edge case just mentioned.

\section{From K215 to consistency in ZFC}\label{sec:logic}
The pinned Lean endpoint is
\begin{lstlisting}
KEndpoint215.blank_halts_iff_zfc_refutation :
  BlankHalts machine215 <-> BProv ZFCax_s [] Form.fBot
\end{lstlisting}
The Rocq development gives the corresponding inconsistency endpoint.
The upstream verification README distinguishes these external
proof-assistant theorems from a proof encoded inside the object's own
ZFC syntax~\cite{k215,endpoint}. Its reported Lean axioms are
\file{propext}, \file{Classical.choice} and \file{Quot.sound}; the
Rocq endpoint reports are closed. These are upstream records, not
fresh builds performed for this packet.

The ordinary mathematical calibration is formalizable in ZFC. Its
numerical derivability predicate is inductively generated over the
natural numbers; the set-theoretic calibration uses first-order
completeness, set models, relative constructibility, reflection and
quotient constructions. The simulation, finite-window compiler,
Pell arguments and coefficient computations are ordinary finite
combinatorics and arithmetic. No proper-class model or inaccessible
universe is needed to transfer this chain to ZFC.

Combining the three equivalences gives
\[
 \exists\boldsymbol{x}\;\mathcal S(\boldsymbol{x})
 \ \longleftrightarrow\
 \mathcal M\text{ accepts raw input }2
 \ \longleftrightarrow\
 K215\text{ halts blank}
 \ \longleftrightarrow\
 \neg\Con(\ZFC).
\]
Each equivalence, including the finite definition of the coefficients,
is provable in ZFC. Negating yields Theorem~\ref{thm:main}.

This is an equivalence theorem, not a proof of consistency.
If ZFC is consistent, the constructed equations have no solutions in
the standard integers, and ZFC cannot prove that lack of solutions by
G\"odel's second incompleteness theorem. Bare consistency alone should
not be advertised as proving independence in both directions:
excluding a ZFC proof of the false $\Sigma_1$ assertion of a solution
requires an additional soundness assumption, such as $\Sigma_1$-soundness.

\section{Exact degree and source counts}\label{sec:counts}
The seven-equation circuit has 34 multiplications and 29 additions or
subtractions. All sixteen unknowns, all five coefficient ports and all
63 gates are live. The exact equation degrees, in the order of
\eqref{eq:system}, are
\[
 22,\ 18,\ 28,\ 7,\ 2,\ 22,\ 1.
\]
The apparent degree-22 leading terms in the main norm cancel; its true
degree is 18. Counts based solely on syntactic degree propagation would
therefore be inaccurate.

Forming the seven residuals, squaring them and adding them uses another
$7+7+6=20$ operations, giving 83 for $P_{\rm SOS}$. Its exact degree
is 56 because the auxiliary residual is the unique one of degree 28.

For $P_{75}$ the factor degrees are
$22,18,28,7,2,22,2$, summing to 101.
Put $Q=(B-1)J$ and $k_0=\eta+\zeta$. Its top homogeneous form is
\begin{align*}
 8Q^{55}h\gamma i^2k_0^9w^7s^{16}T^2f^4
 (wZ-tQ)(Q-F-2Z-\alpha)^2.
\end{align*}
The monomial
$J^{58}h\gamma i^2\eta^9w^7s^{16}tT^2f^4$
has coefficient $-8(B-1)^{58}\ne0$. Thus degree 101 is exact on every
admissible coefficient slice. The circuit ledger is $41M+34A=75$.

For $P_{73}$ the factor degrees are
$22,18,60,7,2,46,2$, summing to 157. Its top homogeneous form is
\[
 -8Q^{91}h\gamma i^4k_0^{13}w^{13}s^{26}T^2f^2
 (wZ-tQ)(Q-F-2Z-\alpha)^2.
\]
The monomial
$J^{94}h\gamma i^4\eta^{13}w^{13}s^{26}tT^2f^2$
has coefficient $8(B-1)^{94}\ne0$. Subtracting $\Delta$, of degree
12, cannot affect it. The ledger is $41M+32A=73$.

The final squared-marker versions delete exactly one addition each.
Their ledgers are therefore $41M+33A=74$ and $41M+31A=72$.
Both marker factors have the same leading homogeneous form $r_0^2$,
so the exact degrees remain 101 and 157. The old 75- and 73-operation
sources are retained as proof parents and arithmetic controls.

The reduction from the universal F84 source removes a variable-input
Pell bridge that is unnecessary once arithmetic input 1 is fixed.
The fixed marker $W_0$ replaces its role, and the positive main
coordinate $\gamma=\rho+\sigma$ replaces two coordinates.
The raw slack restoration in \eqref{eq:rawslack} permits a further
subtraction to disappear. This is why the result is specifically a
fixed predicate, rather than a new ordinary-input universal circuit
with the same counts.

\section{Domains and the explicit quartic alternative}\label{sec:quartic}

\subsection{Nonnegative and unrestricted integer domains}
For a standard nonnegative domain $\N=\{0,1,2,\ldots\}$, replace each
positive primary variable $x_j$ by $1+v_j$, $v_j\in\N$.
This preserves the variable count, polynomial degree and solvability.
The compact operation counts in Table~\ref{tab:tradeoffs} refer to
the original positive-input interfaces; shifts are not silently free
additional circuit instructions.

For an equation over unrestricted integers, Lagrange's four-square
theorem gives a direct explicit conversion. For each $j=1,\ldots,16$
put
\[
 x_j=1+u_{j1}^2+u_{j2}^2+u_{j3}^2+u_{j4}^2,
 \qquad u_{jr}\in\Z.
\]
Substitute these expressions into $P_{\rm SOS}$. The resulting
single ordinary polynomial in 64 unrestricted integer unknowns has
degree 112 and has a zero if and only if $\mathcal S$ has a positive
solution. Equivalently, substitute into the seven equations to get
a system with maximum degree 56. The construction and degrees are
explicit; no positivity predicate remains outside the polynomial.
The complete 211-operation circuit is supplied in
\file{integer64_sos.json}; \file{emit_sos_variants.py} regenerates it
and the 83-operation positive-domain SOS source.

\subsection{A smaller-degree lift over the nonnegative integers}
The circuit of $\mathcal S$ has 63 gates, but a naive lift cannot
automatically make all their values nonnegative: subtraction gates
need not be positive on arbitrary assignments. Here a solution-specific
audit succeeds after one exact sign transformation.

Replace the three auxiliary instructions
\[
 G_a=V^2-y^2,\quad H_a=K_aG_a,\quad N_a=H_a+y^2
\]
by
\[
 G_a=y^2-V^2,\quad H_a=K_aG_a,\quad N_a=y^2-H_a.
\]
The final auxiliary polynomial is identically unchanged, while two
internal registers change sign. Section~\ref{sec:proof} proves $y>V>1$
at every positive solution. The complete audit in
\file{QUADRATIC_LIFT.md} checks every subtraction gate; all 63
modified gate values are nonnegative at such a solution.

Introduce $v_1,\ldots,v_{16}\in\N$ with $x_j=1+v_j$, and one
nonnegative unknown $g_r$ for each gate. An instruction $r=u\circ v$
becomes
\[
 g_r-\widehat u\circ\widehat v=0,
 \qquad\circ\in\{+,-,\cdot\},
\]
where each operand is a fixed coefficient, an affine primary expression
$1+v_j$, or an earlier gate variable. The resulting 63 gate equations
have degree at most two. Add the seven final comparisons. This gives
79 nonnegative unknowns and 70 quadratic equations.

The forward implication uses the proven nonnegativity of each gate
at a solution. Conversely, topological induction on the gate equations
recovers the evaluated circuit and then every equation of
$\mathcal S$. Thus no extra or lost solutions are introduced.

If $E_1,\ldots,E_{70}$ are the explicitly emitted residuals, then
\begin{equation}\label{eq:quartic}
 Q_{79}=\sum_{r=1}^{70}E_r^2=0
\end{equation}
is a single polynomial of degree four with the same solvability
predicate over $\N^{79}$. The first gate equation is
$g_{\rm tau\_square}-(1+v_{\rm tau\_root})^2=0$.
Its square supplies the coefficient
$[v_{\rm tau\_root}^4]Q_{79}=1$, proving exact degree four.
\file{quadratic_lift.json} contains all 70 residuals and the entire
expanded sparse quartic: 286 formal terms before the five fixed
coefficient ports are specialized. The ports remain coefficients,
not extra unknowns.

\section{Reproducibility, evidence and limitations}\label{sec:repro}
The archive contains the complete \LaTeX{} source and PDF, standalone
Python programs using only the standard library, literal arithmetic
circuits, an expanded quartic, exact coefficient definitions, the
K215 table, proof notes, review notes and machine-readable receipts.
The \file{README.md} gives the file map and commands.

The arithmetic validators authenticate the parent F84 JSON by its Git
blob hash, reconstruct the child row changes, check closure and
liveness, recount every operation, compare signed full-source
identities, and verify exact degree by integer polynomial arithmetic.
The ordinary variant also checks the strong-coordinate transformation.
These checks use small diagnostic coefficient tuples only as algebraic
controls; those tuples are not claimed to encode ZFC.

The normalized machine checks all 3,488 rows, 5,160 folded transition
cases, four input erasures, and matched finite simulation prefixes.
The exact window-count formula was checked against exhaustive
five-position enumeration on two small machines. The finite compiler
was compared field by field with the two pinned reference compiler
functions on five fixtures, covering both correction-bit branches.
On the smallest fixture the complete fixed numerals were materialized,
and native mask ranges, populations and modular evaluations checked.

The quartic generator proves the local sign identity by sparse
polynomial arithmetic, emits all residuals, expands their squares,
checks the exact degree and coefficient, and compares the result on
64 signed circuit assignments. Separate reviews rederived the main
sign and index arguments, exact degree leaders and coefficient layout.
These are reviews and tests produced within this work, not a claim
of independent human peer review.

\paragraph{What has not been computed or formalized.}
The actual K215 window list, expanded giant coefficient tuple,
correction bit and any accepting witness have not been materialized.
No test decides K215 halting or ZFC consistency. The new end-to-end
Diophantine theorem has not been compiled in Lean or Rocq, and the
upstream proof assistants were not rerun here. The proof is an ordinary
ZFC-formalizable mathematical proof with the explicit pinned
dependencies listed below. This distinction is material when deciding
where the packet belongs in the repository's research and verification
directories.

\appendix
\section{Source interface and reproducible commands}
\begin{longtable}{@{}ll@{}}
\toprule Mathematical variable & Literal source name\\\midrule
\endhead
$J$ & \file{Jrep}\\
$F$ & \file{F}\\
$\alpha$ & \file{alpha}\\
$t$ & \file{transport_quotient}\\
$f$ & \file{f}\\
$h$ & \file{h}\\
$i$ & \file{i}\\
$T$ & \file{auxiliary_quotient}\\
$s$ & \file{s}\\
$w$ & \file{w}\\
$\tau$ & \file{tau_root}\\
$\eta$ & \file{eta}\\
$\zeta$ & \file{zeta}\\
$y$ & \file{y_aux}\\
$Z$ & \file{Z}\\
$\gamma$ & \file{gamma}\\\bottomrule
\end{longtable}
From the extracted directory, run
\begin{lstlisting}
python3 reproduce.py
latexmk -pdf -interaction=nonstopmode -halt-on-error report.tex
\end{lstlisting}
The reproduction command performs the finite implementation and
arithmetic checks. It never enumerates the actual huge window list or
materializes the actual K215 coefficient tuple. The included JSON
receipts record the scope of each check. The PDF build requires a
normal \TeX{} installation; the Python checks use no third-party packages.

\section{Pinned mathematical dependencies}
Every ProveIt reference below is fixed to commit
\begin{center}\ttfamily cdd582bec13f908df3711fa9e05d1ddac368a5db.\end{center}
The repository base is \repo. The full paths in the bibliography can
be appended to its \file{blob/<commit>/} route. This avoids an evolving
\file{main} branch changing the meaning of the cited theorem.
The supplied \file{PROVENANCE.json} records source locations and hashes;
\file{ATTRIBUTION.md} explains copied notices and inherited terms.

\begin{thebibliography}{99}\footnotesize
\setlength{\itemsep}{2pt}
\setlength{\parskip}{0pt}
\bibitem{k215}
V. Reshetnikov and ProveIt contributors.
\emph{K215 verification status and literal candidate}.
\path{Computability/BusyBeaver/Research/ZFC/k225-formal/verification/k215/README.md};
\path{Computability/BusyBeaver/Research/ZFC/candidate-k215/README.md};
\path{Computability/BusyBeaver/Research/ZFC/candidate-k215/k215.tm}.

\bibitem{endpoint}
ProveIt contributors. \emph{K215 consistency calibration endpoint}.
\path{Computability/BusyBeaver/Lean/BusyBeaver/ZFC/KEndpoint215.lean},
with its imported calibration and machine proofs.

\bibitem{raw75}
ProveIt contributors. \emph{A 75-operation fixed-index universal certificate}.
\path{Computability/HilbertTenthProblem/Papers/1980/FIXED_RAW_UNIVERSAL_75_PROOF.md}.
This is the complete raw arithmetic/compiler theorem used for soundness.

\bibitem{f84}
ProveIt contributors. \emph{Scaled strong output, complete F84 source}.
\path{Computability/HilbertTenthProblem/Papers/research-wip/native-stream-queue/complete84_scaled_strong_output.md}
and the adjacent \file{.json} source. The source JSON Git blob is
\texttt{a8e6e50e96503e69221b8586e273208e5345eee4}.

\bibitem{rank}
ProveIt contributors. \emph{Relaxed auxiliary Pell proof}.
\path{Computability/HilbertTenthProblem/Papers/1980/PELL_RELAXED_AUXILIARY_PROOF.md}.
See also
\path{Computability/HilbertTenthProblem/Papers/research-wip/native-stream-queue/complete75_coupled_index_linear88.md}
for the generic strong-rank application before index recovery.

\bibitem{reverse}
ProveIt contributors. \emph{Reversed auxiliary construction}.
\path{Computability/HilbertTenthProblem/Papers/research-wip/native-stream-queue/complete75_reversed_auxiliary89.md},
especially the signed index step-down in Section 4.

\bibitem{asymmetric}
ProveIt contributors. \emph{Asymmetric scale tradeoffs}.
\path{Computability/HilbertTenthProblem/Papers/research-wip/native-stream-queue/complete75_asymmetric_scale_tradeoffs.md},
especially the input-independent power recovery.

\bibitem{quotient}
ProveIt contributors. \emph{Auxiliary quotient and normalized strong variants}.
\path{Computability/HilbertTenthProblem/Papers/research-wip/native-stream-queue/complete85_auxiliary_bezout_projection.md};
\path{Computability/HilbertTenthProblem/Papers/research-wip/native-stream-queue/complete75_normalized_strong87.md}.

\bibitem{compiler}
ProveIt contributors. \emph{Finite helical window compiler and modified masks}.
\path{Computability/HilbertTenthProblem/Papers/verification/explore_fixed_raw_universal_76.py};
\path{Computability/HilbertTenthProblem/Papers/research-wip/native-stream-queue/complete75_half_binomial_compiler.py}
and its adjacent explanatory \file{.md} file.

\bibitem{nine}
ProveIt contributors. \emph{Nine-unknown representation theorem}.
\path{Computability/HilbertTenthProblem/Lean/Diophantine/Paper1982/NineMain.lean},
theorem \file{Jones1982.nine_unknowns}.

\bibitem{jones}
J. P. Jones. Universal Diophantine equation.
\emph{Journal of Symbolic Logic} 47 (1982), 549--571.
\url{https://www.cambridge.org/core/journals/journal-of-symbolic-logic/article/abs/universal-diophantine-equation/113EBF3D928DFC6011C52BECAE81A087}.

\bibitem{carlmoroz}
M. Carl and B. Z. Moroz. On a Diophantine representation of the predicate
of provability. \emph{Journal of Mathematical Sciences} 199 (2014), 36--52.
\url{https://doi.org/10.1007/s10958-014-1830-2}.

\bibitem{fgo}
A. Formisano, I. Gavazzi and E. G. Omodeo.
Diophantine modeling of provability in algebraic logic.
\emph{CILC 2025}, CEUR Workshop Proceedings 4003.
\url{https://ceur-ws.org/Vol-4003/paper13.pdf}.
These last two papers provide related explicit provability constructions;
their equations are not used as uninstantiated coefficients here.
\end{thebibliography}
\end{document}
